\section{関連研究}\label{sec2}

\subsection{LLMを活用した家庭内タスクプランニング}\label{sec2.1}
近年，エージェントが，人間から与えられた自然言語の指示に従って，タスクを実行するための行動計画をすることを目的とした研究分野において，LLMの活用が注目を集めている．
LLMの活用には以下の複数の利点が挙げられる．
\begin{itemize}
    \item LLMの言語理解能力によって，エージェントは人間から与えられる「言葉」の多様な言い回しに対応できるようになる．これにより，文脈に応じた適切な解釈が可能となる．
    \item LLMの常識知識と推論能力を利用することで，行動計画の候補を生成できる．ただし，環境状態の欠落や誤解釈が残るため，性能向上はタスクと提示条件に依存する．
    \item 従来手法~\cite{zhang-chai-2021-hierarchical,pmlr-v164-blukis22a}と異なるデータ効率の可能性が報告されている~\cite{song2023llm}が，少量データで高性能になることを本研究の結果から一般化しない．
\end{itemize}

Huangら~\cite{huang22}は，GPT-3とBERT~\cite{devlin2018bert}を用いて，自然言語で表現された抽象的なタスク（例：Make breakfast）をシミュレータ上で実行可能なアクションシーケンスに変換する方法を提案している．Huangらの手法では，初めに，Planning LM (GPT-3)を用いて，タスク達成に有効な次のアクションを一つ自由形式で生成する．その後，Translation LM (BERT)を用いて，Planning LMが生成したアクションをシミュレータで実行可能な形式に変換している．二巡目以降は，直前までのアクションも参照してプランニングを\HL{再帰}的に行う．

しかし，タスクを達成するためのアクションを生成して実行する際，生成したアクションが必ずしも実行可能である状況であるとは限らないといった課題が生じる．例えば，リンゴを手に保持していない状態で，``Put the apple on the table''を生成して実行する可能性がある．このようなプランニング中に発生する前提条件エラーを対処するための手法~\cite{raman2024cape}を紹介する．

Ramanら~\cite{raman2024cape}は，Huangらの手法を拡張し，プランニング中のアクション実行時に発生する前提条件エラーに対処するための修正アクションを自動生成する手法を提案している．Ramanらの手法は，シミュレータからのエラーフィードバックとFew-shot Reasoningを活用することで，生成されるアクションの質を向上させている．これにより，ベースライン手法~\cite{ahn2022can,huang22}と比較して，より多くのタスクを正確に達成した．これに対して本研究では，アクションを実行する前の（アクションを生成した）段階でエラーを予測することにより，実行時のエラーを未然に防ぎ，代替アクションを再生成する手法を提案する．

これらの研究~\cite{huang22,raman2024cape}はこの分野においてLLMが優れた能力を持つことを証明したが，いずれも詳細な環境情報を扱っていない．Linら~\cite{lin2023grounded}は，LLMが環境内の認識を行い，その理解に基づいてタスクプランニングをすることに焦点を当てた．Linらは，シミュレータから取得可能な環境内の物体に関する状態や位置などのテーブルデータを用いて，エージェントがタスクを実行するためのアクションを逐次的に計画する手法を提案している．これに対して本研究では，家庭環境に関する具体的な知識を限定し，それらを自然言語に変換してLLMに提供する．


\subsection{家庭内タスクベンチマーク}\label{sec2.2}
人間から与えられる自然言語の指示から，タスクを実行するための適切な行動計画を行う研究のための代表的なデータセットとして，VirtualHome (VH) データセット~\cite{Puig_2018_CVPR}とALFRED~\cite{ALFRED20}の二つがある．

VHデータセットは，エージェントに家庭内での活動を学習させるために収集したデータセットであり，人間が家庭で行う活動（例：コーヒーを入れる）と，その活動を実行することができる様々な方法を自然言語で記述した複数の説明，それぞれの説明に対してエージェントがアクティビティを実行するために必要なアクションシーケンスによって構成されている．VHデータセットを用いた評価では，Longest Common Subsequence (LCS) を用いて正解との一致度を測る指標~\cite{Puig_2018_CVPR}が一般的に用いられる．しかし，実際にエージェントが家事をすることを想定した場合，その時の状況によって同じ指示でも振舞うべき行動が変わる可能性が高いが，これを想定した評価ができない．

ALFREDは，自然言語による指示とエージェントの視覚情報から，家庭内タスクを実行するためのアクションシーケンスへのマッピングを学習するためのベンチマークである．ALFREDデータセットは，デモンストレーションに対応する自然言語の指示で構成されている．各指示にはタスクと，段階的な手順が含まれている．各デモンストレーションは，エージェント視点から観察した視覚情報と，各時間ステップで取ったアクションシーケンスで構成されており，AI2THOR 2.0~\cite{ai2thor}で再生可能なコマンドで記述される．ALFRED では，Task Success と Goal-Condition Success を評価する．Task Success は，物体の位置や状態がタスクのゴール条件に正しく対応していれば１，そうでなければ０と定義される．Goal-Condition Success は，タスク完了に必要なゴール条件に対する完了したゴール条件の割合である．

ただし，これらのデータセットの主目的は本研究と異なり，物体の現在状態・識別番号・関係を明示的に入力して前提条件を検証することに限定されない．本研究はその能力を診断する補完的なタスク設定として位置付け，既存データセットに同様のタスクが存在しないとは断定しない．

\begin{table}[tb]
\caption{関連研究との評価対象の比較（数値の直接比較ではない）}\label{tab:related_comparison}
\centering\small
\begin{tabular}{p{.22\linewidth}p{.22\linewidth}p{.22\linewidth}p{.22\linewidth}}
\hline
研究 & 環境入力 & 実行時修正 & 本研究との差分 \\
\hline
Huangら~\cite{huang22} & 抽象タスクと履歴 & 主に形式変換 & 状態・ID・関係の診断なし \\
Ramanら~\cite{raman2024cape} & シミュレータ応答 & エラー後修正 & 本研究は実行前検証も評価 \\
Linら~\cite{lin2023grounded} & 状態・位置の表形式 & 逐次計画 & 表現形式とタスク集合が異なる \\
本研究 & 状態・関係の自然言語化 & 検証後1回再生成 & 4タスクのログ付き診断を併記 \\
\hline
\end{tabular}
\end{table}
